\documentclass[12pt]{article}
\usepackage{amsmath}
\usepackage{geometry}
\geometry{margin=1in}
\title{Pressure Worksheet}
\date{}
\begin{document}
\maketitle

\section*{Questions}
\begin{enumerate}
\item Starting from the pressure formula \(P = \frac{F}{A}\), rewrite the formula so that:
    \begin{enumerate}
        \item \(A\) is the subject.
        \item \(F\) is the subject.
    \end{enumerate}

\item A force of \(250\,\text{N}\) acts on an area of \(0.50\,\text{m}^2\). Calculate the pressure.

\item A pressure of \(4000\,\text{Pa}\) acts on a surface of area \(2.0\,\text{m}^2\). Calculate the force exerted.

\item A force of \(100\,\text{N}\) produces a pressure of \(2.0\times10^{3}\,\text{Pa}\). Calculate the area.

\item A box weighing \(600\,\text{N}\) rests on the floor. The area of contact is \(0.80\,\text{m}^2\). Calculate the pressure exerted on the floor.

\item A person of weight \(650\,\text{N}\) is standing on a pair of snowshoes with a total area of \(0.30\,\text{m}^2\). Calculate the pressure on the snow.

\item A hydraulic press applies a pressure of \(3.0\,\text{kPa}\) to a metal plate of area \(0.25\,\text{m}^2\). Calculate the force on the plate.

\item A brick exerts a pressure of \(5200\,\text{Pa}\) on the ground. If the weight of the brick is \(130\,\text{N}\), calculate the area of contact with the ground.

\item A force of \(2.5\,\text{kN}\) is applied uniformly to an area of \(0.40\,\text{m}^2\). Calculate the pressure in \(\text{kPa}\).

\item A piston is subjected to a pressure of \(850\,\text{kPa}\). The piston is circular with a diameter of \(10\,\text{cm}\). (Use \(A = \pi r^{2}\).) Calculate the force on the piston.

\item A table weighs \(1200\,\text{N}\) and stands on four rectangular feet, each measuring \(5\,\text{cm} \times 7\,\text{cm}\). (Use \(A = \text{width} \times \text{length}\).) Calculate the pressure exerted by one foot on the floor. (Assume the weight is shared equally between the four feet.)

\item The circular base of a post has radius \(12\,\text{cm}\) and exerts a pressure of \(15\,\text{kPa}\) on the ground. (Use \(A = \pi r^{2}\).) Calculate the force exerted by the base.
\end{enumerate}

\newpage
\section*{Answer Key}
\begin{enumerate}
\item \(A = \dfrac{F}{P}\);\qquad \(F = P A\)
\item \(5.0\times10^{2}\,\text{Pa}\)
\item \(8.0\times10^{3}\,\text{N}\)
\item \(5.0\times10^{-2}\,\text{m}^2\)
\item \(7.5\times10^{2}\,\text{Pa}\)
\item \(2.2\times10^{3}\,\text{Pa}\;(2.2\,\text{kPa})\)
\item \(7.5\times10^{2}\,\text{N}\)
\item \(2.5\times10^{-2}\,\text{m}^2\)
\item \(6.25\,\text{kPa}\)
\item \(6.7\times10^{3}\,\text{N}\;(6.7\,\text{kN})\)
\item \(8.6\times10^{4}\,\text{Pa}\;(86\,\text{kPa})\)
\item \(6.8\times10^{2}\,\text{N}\)
\end{enumerate}

\newpage
\section*{Worked Solutions}
\begin{enumerate}
\item From \(P = \dfrac{F}{A}\):
    \begin{align*}
    A &= \frac{F}{P},\\[4pt]
    F &= P A.
    \end{align*}

\item \(P = \dfrac{F}{A} = \dfrac{250\,\text{N}}{0.50\,\text{m}^2} = 5.0 \times 10^{2}\,\text{Pa}.\)

\item \(F = P A = (4000\text{ Pa})(2.0\,\text{m}^2) = 8.0 \times 10^{3}\,\text{N}.\)

\item Rearranging: \(A = \dfrac{F}{P} = \dfrac{100\,\text{N}}{2.0 \times 10^{3}\,\text{Pa}} = 5.0 \times 10^{-2}\,\text{m}^2.\)

\item \(P = \dfrac{600\,\text{N}}{0.80\,\text{m}^2} = 7.5 \times 10^{2}\,\text{Pa}.\)

\item \(P = \dfrac{650\,\text{N}}{0.30\,\text{m}^2} = 2.17 \times 10^{3}\,\text{Pa} \approx 2.2\,\text{kPa}.\)

\item Convert \(3.0\,\text{kPa} = 3.0 \times 10^{3}\,\text{Pa}\).\\
\(F = P A = (3.0 \times 10^{3}\,\text{Pa})(0.25\,\text{m}^2) = 7.5 \times 10^{2}\,\text{N}.\)

\item \(A = \dfrac{F}{P} = \dfrac{130\,\text{N}}{5200\,\text{Pa}} = 2.5 \times 10^{-2}\,\text{m}^2.\)

\item Convert force: \(2.5\,\text{kN} = 2.5 \times 10^{3}\,\text{N}\).\\
\(P = \dfrac{F}{A} = \dfrac{2.5 \times 10^{3}\,\text{N}}{0.40\,\text{m}^2} = 6.25 \times 10^{3}\,\text{Pa} = 6.25\,\text{kPa}.\)

\item Diameter \(10\,\text{cm} \Rightarrow r = 5.0\,\text{cm} = 0.050\,\text{m}\).\\
\(A = \pi r^{2} = \pi(0.050\,\text{m})^{2} = 7.85 \times 10^{-3}\,\text{m}^2.\)\\
Convert pressure: \(850\,\text{kPa} = 850 \times 10^{3}\,\text{Pa}\).\\
\(F = P A = (0.85 \times 10^{6}\,\text{Pa})(7.85 \times 10^{-3}\,\text{m}^2) \approx 6.7 \times 10^{3}\,\text{N}.\)

\item Foot area: \(A = 0.05\,\text{m} \times 0.07\,\text{m} = 3.5 \times 10^{-3}\,\text{m}^2.\)\\
Weight per foot: \(F = \dfrac{1200\,\text{N}}{4} = 300\,\text{N}.\)\\
\(P = \dfrac{300\,\text{N}}{3.5 \times 10^{-3}\,\text{m}^2} \approx 8.57 \times 10^{4}\,\text{Pa} = 86\,\text{kPa}.\)

\item Radius \(r = 12\,\text{cm} = 0.12\,\text{m}\).\\
\(A = \pi r^{2} = \pi(0.12\,\text{m})^{2} = 4.52 \times 10^{-2}\,\text{m}^2.\)\\
Convert pressure: \(15\,\text{kPa} = 1.5 \times 10^{4}\,\text{Pa}\).\\
\(F = P A = (1.5 \times 10^{4}\,\text{Pa})(4.52 \times 10^{-2}\,\text{m}^2) \approx 6.8 \times 10^{2}\,\text{N}.\)
\end{enumerate}

\end{document}
